% GENERATED FILE. Edit canonical lessons, then run npm run generate:books.
% canonical-source-hash: fnv1a64:000877674c13ee1e
% canonical-lessons: PA-C335-sona, PA-C335-suniar, PA-C335-loha, PA-C335-gharna, PA-C335-misatri

\chapter{Gold and Iron}
\label{ch:pa-a2-gold-and-iron}
\hlchaptermodality{335}

\begin{chapteropening}
\textbf{By the end of this chapter:} \emph{I can name the people who work with metal.}

Complete the last lesson of chapter 335: I can name the people who work with metal.
\end{chapteropening}

\section[\texorpdfstring{\pa{ਸੋਨਾ} (sonā)}{ਸੋਨਾ (sonā)}]{\pa{ਸੋਨਾ} (sonā) — gold}

\label{lesson:PA-C335-sona}



\begin{quote}
Before the new one: say the Punjabi for a weaver, then the Punjabi for to dye.
\end{quote}

\subsection*{You'll want to know: \pa{ਸੋਨਾ}}
\textbf{\pa{ਸੋਨਾ}} — \emph{sonā} — \textquotedblleft{}gold\textquotedblright{}.

\begin{quote}
\pa{ਵਿਆਹ} \pa{ਤੋਂ} \pa{ਪਹਿਲਾਂ} \pa{ਅਸੀਂ} \pa{ਸੋਨਾ} \pa{ਖ਼ਰੀਦਿਆ।}
\end{quote}

\begin{quote}
\emph{viāh ton pahilāṁ asīṁ sonā kharīdiā.}
\end{quote}

\begin{quote}
Before the wedding we bought gold.
\end{quote}

A masculine noun.

\begin{cousinweb}
\textbf{\pa{ਸੋਨਾ}} (\emph{sonā}) is inherited from Sanskrit \emph{suvarṇa}, \textquotedblleft{}golden; gold\textquotedblright{}; CDIAL lists Punjabi \emph{sonā}, \textquotedblleft{}gold\textquotedblright{}, under it (CDIAL 13519).
\end{cousinweb}

\begin{grammarlens}[title={What you've built}]
One more word for this chapter.
\end{grammarlens}

\subsection*{Your turn}
\begin{itemize}
\raggedright
  \item \emph{Say it:} \emph{sonā}
  \item \emph{Say it:} \emph{sonā}, once more
  \item \emph{Say it:} \emph{viāh ton pahilāṁ asīṁ sonā kharīdiā.} — Before the wedding we bought gold.
  \item \emph{Recall:} say the Punjabi for a weaver, then the Punjabi for to dye, then say \emph{sonā} again
\end{itemize}

\begin{tcolorbox}[breakable,colback=teal!4,colframe=teal!35!black,title={Before you move on}]
What does \pa{ਸੋਨਾ} mean? (\textquotedblleft{}Gold\textquotedblright{}.) Say it once more.
\end{tcolorbox}
\section[\texorpdfstring{\pa{ਸੁਨਿਆਰ} (suniār)}{ਸੁਨਿਆਰ (suniār)}]{\pa{ਸੁਨਿਆਰ} (suniār) — a goldsmith}

\label{lesson:PA-C335-suniar}



\begin{quote}
Before the new one: say the Punjabi for to dye, then the Punjabi for gold.
\end{quote}

\subsection*{You'll want to know: \pa{ਸੁਨਿਆਰ}}
\textbf{\pa{ਸੁਨਿਆਰ}} — \emph{suniār} — \textquotedblleft{}a goldsmith\textquotedblright{}.

\begin{quote}
\pa{ਸੁਨਿਆਰ} \pa{ਸੋਨੇ} \pa{ਦੀ} \pa{ਮੁੰਦਰੀ} \pa{ਬਣਾਉਂਦਾ} \pa{ਹੈ।}
\end{quote}

\begin{quote}
\emph{suniār sone dī mundrī baṇāundā hai.}
\end{quote}

\begin{quote}
The goldsmith makes a gold ring.
\end{quote}

A masculine noun.

\begin{cousinweb}
\textbf{\pa{ਸੁਨਿਆਰ}} (\emph{suniār}) is inherited from Sanskrit \emph{suvarṇakāra}, \textquotedblleft{}a goldsmith\textquotedblright{}; CDIAL lists Punjabi \emph{suneār}, \textquotedblleft{}goldsmith\textquotedblright{}, under it (CDIAL 13520).
\end{cousinweb}

\begin{grammarlens}[title={What you've built}]
One more word for this chapter.
\end{grammarlens}

\subsection*{Your turn}
\begin{itemize}
\raggedright
  \item \emph{Say it:} \emph{suniār}
  \item \emph{Say it:} \emph{suniār}, once more
  \item \emph{Say it:} \emph{suniār sone dī mundrī baṇāundā hai.} — The goldsmith makes a gold ring.
  \item \emph{Recall:} say the Punjabi for to dye, then the Punjabi for gold, then say \emph{suniār} again
  \item \emph{Write it:} copy \textbf{\pa{ਸੁਨਿਆਰ}} once, checking the whole word
\end{itemize}

\begin{tcolorbox}[breakable,colback=teal!4,colframe=teal!35!black,title={Before you move on}]
What does \pa{ਸੁਨਿਆਰ} mean? (\textquotedblleft{}A goldsmith\textquotedblright{}.) Say it once more.
\end{tcolorbox}
\section[\texorpdfstring{\pa{ਲੋਹਾ} (lohā)}{ਲੋਹਾ (lohā)}]{\pa{ਲੋਹਾ} (lohā) — iron}

\label{lesson:PA-C335-loha}



\begin{quote}
Before the new one: say the Punjabi for gold, then the Punjabi for a goldsmith.
\end{quote}

\subsection*{You'll want to know: \pa{ਲੋਹਾ}}
\textbf{\pa{ਲੋਹਾ}} — \emph{lohā} — \textquotedblleft{}iron\textquotedblright{}.

\begin{quote}
\pa{ਇਹ} \pa{ਪੁਲ} \pa{ਲੋਹੇ} \pa{ਦਾ} \pa{ਹੈ।}
\end{quote}

\begin{quote}
\emph{eh pul lohe dā hai.}
\end{quote}

\begin{quote}
This bridge is made of iron.
\end{quote}

A masculine noun.

\begin{cousinweb}
\textbf{\pa{ਲੋਹਾ}} (\emph{lohā}) is inherited from Sanskrit \emph{loha}, \textquotedblleft{}red; copper; iron\textquotedblright{}; CDIAL lists Punjabi \emph{lohā}, \textquotedblleft{}iron\textquotedblright{}, under it (CDIAL 11158).
\end{cousinweb}

\begin{grammarlens}[title={What you've built}]
One more word for this chapter.
\end{grammarlens}

\subsection*{Your turn}
\begin{itemize}
\raggedright
  \item \emph{Say it:} \emph{lohā}
  \item \emph{Say it:} \emph{lohā}, once more
  \item \emph{Say it:} \emph{eh pul lohe dā hai.} — This bridge is made of iron.
  \item \emph{Recall:} say the Punjabi for gold, then the Punjabi for a goldsmith, then say \emph{lohā} again
  \item \emph{Write it:} copy \textbf{\pa{ਲੋਹਾ}} once, checking the whole word
\end{itemize}

\begin{tcolorbox}[breakable,colback=teal!4,colframe=teal!35!black,title={Before you move on}]
What does \pa{ਲੋਹਾ} mean? (\textquotedblleft{}Iron\textquotedblright{}.) Say it once more.
\end{tcolorbox}
\section[\texorpdfstring{\pa{ਘੜਨਾ} (ghaṛnā)}{ਘੜਨਾ (ghaṛnā)}]{\pa{ਘੜਨਾ} (ghaṛnā) — to hammer into shape}

\label{lesson:PA-C335-gharna}



\begin{quote}
Before the new one: say the Punjabi for a goldsmith, then the Punjabi for iron.
\end{quote}

\subsection*{You'll want to know: \pa{ਘੜਨਾ}}
\textbf{\pa{ਘੜਨਾ}} — \emph{ghaṛnā} — \textquotedblleft{}to hammer into shape\textquotedblright{}.

\begin{quote}
\pa{ਲੁਹਾਰ} \pa{ਲੋਹੇ} \pa{ਤੋਂ} \pa{ਸੰਦ} \pa{ਘੜਦਾ} \pa{ਹੈ।}
\end{quote}

\begin{quote}
\emph{luhār lohe ton sand ghaṛdā hai.}
\end{quote}

\begin{quote}
The blacksmith hammers tools out of iron.
\end{quote}

Past (I …): \pa{ਮੈਂ} \pa{ਇੱਕ} \pa{ਸੰਦ} \pa{ਘੜਿਆ}, \emph{maiṁ ikk sand ghaṛiā}.

\begin{cousinweb}
\textbf{\pa{ਘੜਨਾ}} (\emph{ghaṛnā}) is inherited from Sanskrit \emph{ghaṭate}, \textquotedblleft{}is busy with\textquotedblright{}; CDIAL lists Punjabi \emph{ghaṛnā}, \textquotedblleft{}to form by hammering\textquotedblright{}, under it (CDIAL 4407).
\end{cousinweb}

\begin{grammarlens}[title={What you've built}]
One more word for this chapter.
\end{grammarlens}

\subsection*{Your turn}
\begin{itemize}
\raggedright
  \item \emph{Say it:} \emph{ghaṛnā}
  \item \emph{Say it:} \emph{ghaṛnā}, once more
  \item \emph{Say it:} \emph{luhār lohe ton sand ghaṛdā hai.} — The blacksmith hammers tools out of iron.
  \item \emph{Recall:} say the Punjabi for a goldsmith, then the Punjabi for iron, then say \emph{ghaṛnā} again
  \item \emph{Write it:} copy \textbf{\pa{ਘੜਨਾ}} once, checking the whole word
\end{itemize}

\begin{tcolorbox}[breakable,colback=teal!4,colframe=teal!35!black,title={Before you move on}]
What does \pa{ਘੜਨਾ} mean? (\textquotedblleft{}To hammer into shape\textquotedblright{}.) Say it once more.
\end{tcolorbox}
\section[\texorpdfstring{\pa{ਮਿਸਤਰੀ} (misatrī)}{ਮਿਸਤਰੀ (misatrī)}]{\pa{ਮਿਸਤਰੀ} (misatrī) — a mechanic}

\label{lesson:PA-C335-misatri}



\begin{quote}
Before the new one: say the Punjabi for iron, then the Punjabi for to hammer into shape.
\end{quote}

\subsection*{You'll want to know: \pa{ਮਿਸਤਰੀ}}
\textbf{\pa{ਮਿਸਤਰੀ}} — \emph{misatrī} — \textquotedblleft{}a mechanic\textquotedblright{}.

\begin{quote}
\pa{ਮਿਸਤਰੀ} \pa{ਸਾਡੀ} \pa{ਗੱਡੀ} \pa{ਠੀਕ} \pa{ਕਰਦਾ} \pa{ਹੈ।}
\end{quote}

\begin{quote}
\emph{misatrī sāḍī gaḍḍī ṭhīk kardā hai.}
\end{quote}

\begin{quote}
The mechanic fixes our car.
\end{quote}

A masculine noun.

\begin{cousinweb}
\textbf{\pa{ਮਿਸਤਰੀ}} (\emph{misatrī}) is a loan, usually traced to Portuguese \emph{mestre}, \textquotedblleft{}a master of a trade\textquotedblright{}.
\end{cousinweb}

\begin{grammarlens}[title={What you've built}]
One more word for this chapter.
\end{grammarlens}

\subsection*{Your turn}
\begin{itemize}
\raggedright
  \item \emph{Say it:} \emph{misatrī}
  \item \emph{Say it:} \emph{misatrī}, once more
  \item \emph{Say it:} \emph{misatrī sāḍī gaḍḍī ṭhīk kardā hai.} — The mechanic fixes our car.
  \item \emph{Recall:} say the Punjabi for iron, then the Punjabi for to hammer into shape, then say \emph{misatrī} again
  \item \emph{Write it:} copy \textbf{\pa{ਮਿਸਤਰੀ}} once, checking the whole word
\end{itemize}

\begin{tcolorbox}[breakable,colback=teal!4,colframe=teal!35!black,title={Before you move on}]
What does \pa{ਮਿਸਤਰੀ} mean? (\textquotedblleft{}A mechanic\textquotedblright{}.) Say it once more.
\end{tcolorbox}
